This thesis explored local \gls{llm} assistance for \gls{cli} workflows under
tight resource constraints.
Command-vs.\ natural-language classification was not
reliable enough for real-time use: across \texttt{classifier-100}, accuracy
ranged from 50\% to 88\%, and the best model required 2.9s per query.
Document retrieval with embedding \glspl{llm} over local man pages was more
successful. Across the evaluated corpora, every retrieval strategy achieved good
accuracy, around 75\% to 94\% Top-5 accuracy and 0.6 to 0.9 MRR score on
\texttt{sample-100}. The end-to-end latency testing showed that document-level
retrieval produced the highest latency due to large output size, while
passage-level retrieval and re-ranking reduced latency significantly. Overall,
passage-level retrieval gave the best balance of latency and retrieval quality,
when using bigger embedding models, while document-level retrieval with
re-ranking was more effective for smaller embedding models, the vector
computation cost for which was much lower. On \texttt{sample-1000}, best setups
reusing models tested on \texttt{sample-100} managed to keep latencies for
search under 2s and RAG under 10s.
The results show that to be practical, local \gls{llm} assistance must be
constrained to specific tasks like retrieval and summarization, rather than
open-ended classification or reasoning.
